\newcommand{\normal}{\fontsize{12pt}{16pt}\selectfont}
\newcommand{\size}{\fontsize{14pt}{18pt}\selectfont}
\newcommand{\bigsize}{\fontsize{16pt}{20pt}\selectfont}
\newcommand{\bigbigsize}{\fontsize{20pt}{24pt}\selectfont}
\begin{flushleft}
\size\textbf{PROGRAM OUTCOMES (POs)}
\end{flushleft}
\\
Engineering graduates will be able to:
\vspace{5mm}\\
\textbf{PO1: Engineering knowledge:} Apply knowledge of mathematics, science, engineering fundamentals, and an engineering specialization to solve complex engineering problems.\\
\textbf{PO2: Problem analysis:} Identify, formulate, review research literature, and analyze complex engineering problems to reach substantiated conclusions using mathematics, natural sciences, and engineering sciences.\\
\textbf{PO3: Design/development of solutions:} Design solutions for complex engineering problems and develop system components or processes that meet specified needs with consideration for public health, safety, cultural, societal, and environmental concerns.\\
\textbf{PO4: Conduct investigations of complex problems:} Use research-based knowledge and research methods, including design of experiments, analysis and interpretation of data, and synthesis of information, to provide valid conclusions.\\
\textbf{PO5: Modern tool usage:} Create, select, and apply appropriate techniques, resources, and modern engineering and IT tools, including prediction and modeling, to complex engineering activities with an understanding of their limitations.\\
\textbf{PO6: The engineer and society:} Apply contextual knowledge to assess societal, health, safety, legal, and cultural issues and the responsibilities relevant to professional engineering practice.\\
\textbf{PO7: Environment and sustainability:} Understand the impact of professional engineering solutions in societal and environmental contexts and demonstrate the need for sustainable development.\\
\textbf{PO8: Ethics:} Apply ethical principles and commit to professional ethics, responsibilities, and norms of engineering practice.\\
\textbf{PO9: Individual and teamwork:} Function effectively as an individual, and as a member or leader in diverse teams and multidisciplinary settings.\\
\textbf{PO10: Communication:} Communicate effectively on complex engineering activities with the engineering community and society at large, including writing reports, preparing design documentation, making presentations, and giving and receiving clear instructions.\\
\textbf{PO11: Project management and finance:} Demonstrate knowledge of engineering and management principles and apply them to one's own work, as a member or leader in a team, to manage projects in multidisciplinary environments.\\
\textbf{PO12: Life-long learning:} Recognize the need for independent and lifelong learning and develop the ability to pursue it in the broad context of technological change.
\\
\\
\textbf{PROGRAM EDUCATIONAL OUTCOMES (PEOs)}
\\
\\
\textbf{PEO1:} Apply computational skills needed to analyze, formulate, and solve engineering problems.
\\
\textbf{PEO2:} Establish entrepreneurial ability and work in interdisciplinary research and development organizations, individually or as part of a team.
\\
\textbf{PEO3:} Inculcate ethical values and leadership qualities needed for a successful career.
\\
\textbf{PEO4:} Develop analytical thinking for understanding and solving real-world problems and sustain an attitude of lifelong learning for higher education.
\\
\\
\textbf{PROGRAM SPECIFIC OUTCOMES (PSOs)}
\\
\\
\textbf{PSO1:} Acquire in-depth knowledge of theoretical foundations and issues in Computer Science to build learning ability and computational skills.
\\
\textbf{PSO2:} Analyze, design, develop, test, and manage complex software systems and applications using advanced tools and techniques.
